\documentclass[12pt,twoside]{article}
\usepackage{chenart}

% \doublespacing
\onehalfspacing



\title{Simple Article Template}
\author{Texer Chen}
\date{\today}



\begin{document}


\maketitle






\begin{abstract}
An abstract is a brief summary of a research article, thesis, review, conference proceeding, or any in-depth analysis of a particular subject and is often used to help the reader quickly ascertain the paper's purpose. When used, an abstract always appears at the beginning of a manuscript or typescript, acting as the point-of-entry for any given academic paper or patent application. Abstracting and indexing services for various academic disciplines are aimed at compiling a body of literature for that particular subject.

\textbf{Keywords:} 
key1; key2
\end{abstract}



\section{English}

English is a language that started in Anglo-Saxon England. It is originally from Anglo-Frisian and Old Saxon dialects. English is now used as a global language. There are about 375 million native speakers (people who use it as their first language) in the world.

Scots is the language closest to English. The vocabulary of English was influenced by other Germanic languages in the early Middle Ages and later by Romance languages, especially French.


\section{中文}

汉语又称中文或华语，概指由上古汉语（先秦雅言）发展而来、书面使用汉字的分析语，为汉藏语系最大的一支语族。如把整个汉语族视为单一语言，则汉语为世界使用人数最多的语言，目前全世界有五分之一人口将汉语做为母语或第二语言。

汉语在以其做为母语的地方有不同通称，且有多种方言变体，其中以北方汉语为基础的官话最为流行，其衍生而来的现代标准汉语（有国语、普通话、新加坡标准华语等变体）是大中华区的主要通用语。此外，汉语还是联合国正式语文，并被上海合作组织等国际组织采用为官方语言。国际上常将“汉语”称为“中文”，这是因为大多数语言使用表音文字，对于“文”与“语”并不作区分，不符合汉语语法；然而华人使用语素文字，文字并不等于语言，作为母语人士不宜随之误用。


\section{Equations}
\begin{equation}
    \frac{x^{2}}{a^{2}}-\frac{y^{2}}{b^{2}}=1 
\end{equation}
\begin{equation}
    \int_{0}^{1} k\mathrm{d}x = kx+C 
\end{equation}
\begin{equation}
    x ={-b \pm \sqrt{b^2-4ac}\over 2a} 
\end{equation}



\section{Figures}


\begin{figure}[htbp]
    \centering
    \includegraphics[width=.4\linewidth]{example-image}
    \caption{Caption}
    \label{fig:enter-label}
\end{figure}

\begin{figure}[htbp]
    \centering
    \subfigure[sub caption 1]{
    \includegraphics[width=.25\linewidth]{example-image-a}
    }
    \subfigure[sub caption 2]{
    \includegraphics[width=.25\linewidth]{example-image-b}
    }
    \subfigure[sub caption 3]{
    \includegraphics[width=.25\linewidth]{example-image-c}
    }
    \caption{Caption}
    \label{fig:enter-label}
\end{figure}


\section{Tables}

\begin{table}[H]
    \centering
    \begin{tabular}{c|c}
    \toprule
         &  \\
    \midrule
         & \\
    \bottomrule
    \end{tabular}
    \caption{Caption}
    \label{tab:my_label}
\end{table}






\end{document}
